\begin{abstract}




Two major sources of training data exist for post-training modern language models: online~(model-generated rollouts) data, and offline~(human or other-model demonstrations) data. 
These two types of data are typically used by approaches like Reinforcement Learning (RL) and Supervised Fine-Tuning (SFT), respectively.
In this paper, we show that these approaches are not in contradiction, but are instances of a single optimization process.
We derive a \framework, and present the calculations of a wide spectrum of post-training approaches as the gradient of a common objective under different data distribution assumptions and various bias-variance tradeoffs. The gradient estimator is constructed with four interchangeable parts: stabilization mask, reference policy denominator, advantage estimate, and likelihood gradient. 
Motivated by our theoretical findings, we propose \method~(HPT), an algorithm that dynamically selects different training signals.
\mot is designed to yield both effective exploitation of demonstration and stable exploration without sacrificing learned reasoning patterns.
We provide extensive experiments and ablation studies to verify the effectiveness of our unified theoretical framework and \mot.
Across six mathematical reasoning benchmarks and two out-of-distribution suites, \mot consistently surpasses strong baselines across models of varying scales and families.





\end{abstract}
